\section{Randomized --- BPLIn PWords}
%%% generated by the blueprint pipeline from TCSlib.Complexity.Randomized.BPLInPWords
\subsection*{Randomized --- BPLIn PWords: Introduction}

Declarations of the Lean module \texttt{TCSlib.Complexity.Randomized.BPLInPWords}.

\subsection*{Randomized --- BPLIn PWords: Results}

\begin{definition}[Encoded sum step]
\lean{Randomized.BPLInP.sumStep}
\label{Randomized.BPLInP.sumStep}
\leanok
\uses{Complexity.BinArith.add, Complexity.pairFstD, Complexity.pairSndD, Turing.pairEncode}
For a Boolean string $S$, define one sum step by extracting two components from the
nested paired data in $S$, adding their encodings, and pairing the result with the
remaining extracted component.
\end{definition}

\begin{definition}[Iterated encoded sum]
\lean{Randomized.BPLInP.sumW}
\label{Randomized.BPLInP.sumW}
\leanok
\uses{Complexity.pairFstD, Randomized.BPLInP.sumStep, Turing.pairEncode}
For a Boolean string $w$, define its encoded sum by starting with the pairing of the
empty string and $w$, applying the sum-step operation $(|w|+1)$ times, and extracting
the first component of the resulting paired string.
\end{definition}

\begin{theorem}[Iterated Sum Update]
\lean{Randomized.BPLInP.iterate_sumStep}
\label{Randomized.BPLInP.iterate_sumStep}
\leanok
\uses{Complexity.BinArith.add, Complexity.LogProg.bitsVal, Complexity.LogProg.bitsVal_bits, Complexity.pairFstD, Complexity.pairSndD, Randomized.BPLInP.sumStep, Turing.pairEncode}
\difficulty{4}
For every $a,i\in\mathbb N$ and bit string $w$, applying the sum update $i$ times to the
pairing of the binary representation of $a$ with $w$ yields the pairing of the binary
representation of $a+\sum_{j=0}^{i-1}\mathrm{val}(\pi_1(\pi_2^{[j]}(w)))$ with
$\pi_2^{[i]}(w)$. Here $\pi_1$ and $\pi_2$ are the total projections from the pair
parser, and the exponent denotes iteration of a function.
\end{theorem}

\begin{theorem}[Summing an Encoded List]
\lean{Randomized.BPLInP.iterate_sumStep_encWords}
\label{Randomized.BPLInP.iterate_sumStep_encWords}
\leanok
\uses{Complexity.BinArith.add, Complexity.BinArith.encWords, Complexity.BinArith.encWords_cons, Complexity.LogProg.bitsVal_bits, Randomized.BPLInP.sumStep, Turing.pairEncode}
\difficulty{4}
For every finite list $L$ of natural numbers and every $a\in\mathbb N$, applying the sum
update $|L|$ times to the pairing of the binary representation of $a$ with the encoded
list of binary representations of the entries of $L$ yields the pairing of the binary
representation of $a+\sum_{n\in L}n$ with the empty bit string.
\end{theorem}

\begin{theorem}[Length of encoded words]
\lean{Randomized.BPLInP.two_mul_length_le_encWords}
\label{Randomized.BPLInP.two_mul_length_le_encWords}
\leanok
\uses{Complexity.BinArith.encWords, Complexity.BinArith.encWords_cons, Turing.length_pairEncode}
\difficulty{4}
For every finite list $L$ of finite Boolean words, the length of its encoded
concatenation is at least twice the number of words in $L$.
\end{theorem}

\begin{theorem}[Encoded words represent a sum]
\lean{Randomized.BPLInP.sumW_encWords}
\label{Randomized.BPLInP.sumW_encWords}
\leanok
\uses{Complexity.BinArith.add, Complexity.BinArith.encWords, Complexity.LogProg.bitsVal, Complexity.LogProg.bitsVal_bits, Complexity.pairFstD_nil, Complexity.pairFstD_pairEncode, Complexity.pairSndD, Randomized.BPLInP.iterate_sumStep_encWords, Randomized.BPLInP.sumStep, Randomized.BPLInP.sumW, Randomized.BPLInP.two_mul_length_le_encWords, Turing.pairEncode}
\difficulty{5}
For every finite list $L$ of natural numbers, encode each entry $m$ by its binary bit
string and concatenate the resulting self-delimiting word encodings. The operation that
sums the encoded words and decodes the resulting value returns the binary representation
of $\sum_{m\in L}m$.
\end{theorem}

\begin{theorem}[Polynomial-Time Computability of the Summation Function]
\lean{Randomized.BPLInP.polyTimeComputable_sumW}
\label{Randomized.BPLInP.polyTimeComputable_sumW}
\leanok
\uses{Complexity.BinArith.length_iterate_pairSndD_le, Complexity.BinArith.polyTimeComputable_add, Complexity.BinArith.polyTimeComputable_iterate, Complexity.LogProg.bitsVal, Complexity.LogProg.bitsVal_lt_two_pow, Complexity.LogProg.length_bits_le, Complexity.PolyTimeComputable, Complexity.PolyTimeComputable.comp, Complexity.PolyTimeComputable.pairEncode, Complexity.length_pairFstD_le, Complexity.pairFstD, Complexity.pairSndD, Complexity.polyTimeComputable_const, Complexity.polyTimeComputable_id, Complexity.polyTimeComputable_pairFstD, Complexity.polyTimeComputable_pairSndD, Randomized.BPLInP.iterate_sumStep, Randomized.BPLInP.sumStep, Randomized.BPLInP.sumW, Turing.length_pairEncode, Turing.pairEncode}
\difficulty{6}
The function that sums the two encoded natural numbers in a bit string and returns the
resulting encoding is polynomial-time computable.
\end{theorem}

\begin{definition}[One coding step]
\lean{Randomized.BPLInP.codeStep}
\label{Randomized.BPLInP.codeStep}
\leanok
\uses{Complexity.BinArith.add, Complexity.pairFstD, Complexity.pairSndD, Turing.pairEncode}
For a Boolean word $S$, write $a$ and $b$ for the two components decoded from $S$. The
coding step encodes the pair whose first component is the sum of the encoding of
$(a,[	ext{true}])$ and whose second component is $b$ followed by the encoding of
$(a,[])$.
\end{definition}

\begin{theorem}[Iterated Code Step]
\lean{Randomized.BPLInP.iterate_codeStep}
\label{Randomized.BPLInP.iterate_codeStep}
\leanok
\uses{Complexity.BinArith.add, Complexity.BinArith.encWords, Complexity.BinArith.encWords_append_singleton, Complexity.LogProg.bitsVal, Complexity.LogProg.bitsVal_bits, Randomized.BPLInP.codeStep, Turing.pairEncode}
\difficulty{4}
For every $i\in\mathbb N$, applying the code update $i$ times to the pairing of two
empty bit strings gives the pairing of the binary representation of $i$ with the encoded
list of binary representations of $0,1,\ldots,i-1$.
\end{theorem}

\begin{definition}[Iterated coding word]
\lean{Randomized.BPLInP.codesW}
\label{Randomized.BPLInP.codesW}
\leanok
\uses{Complexity.pairSndD, Randomized.BPLInP.codeStep, Turing.pairEncode}
For natural numbers $B,e$ and a Boolean word $x$, let $w$ be the second component of the
decoding of the result of applying the coding step $B(|x|+1)^e$ times to the encoding of
the pair of empty words.
\end{definition}

\begin{theorem}[Word of successive binary codes]
\lean{Randomized.BPLInP.codesW_eq}
\label{Randomized.BPLInP.codesW_eq}
\leanok
\uses{Complexity.BinArith.encWords, Complexity.pairSndD_pairEncode, Randomized.BPLInP.codesW, Randomized.BPLInP.iterate_codeStep}
\difficulty{1}
For natural numbers $B,e$ and a Boolean list $x$, the word produced by the code
iteration is the concatenation, in order, of the binary expansions of the natural
numbers $j$ with $0\le j<B(|x|+1)^e$.
\end{theorem}

\begin{theorem}[Length Bound for Encoded Range]
\lean{Randomized.BPLInP.length_encWords_range_le}
\label{Randomized.BPLInP.length_encWords_range_le}
\leanok
\uses{Complexity.BinArith.encWords, Complexity.BinArith.length_encWords_le, Turing.length_bits_le_self}
\difficulty{4}
For every $V\in\mathbb N$, the encoding of the list of binary representations of the
integers from $0$ through $V-1$ has length at most $V(2V+2)$.
\end{theorem}

\begin{theorem}[Polynomial-time computability of code generation]
\lean{Randomized.BPLInP.polyTimeComputable_codesW}
\label{Randomized.BPLInP.polyTimeComputable_codesW}
\leanok
\uses{Complexity.BinArith.polyTimeComputable_add, Complexity.BinArith.polyTimeComputable_iterate, Complexity.PolyTimeComputable, Complexity.PolyTimeComputable.append, Complexity.PolyTimeComputable.comp, Complexity.PolyTimeComputable.pairEncode, Complexity.polyTimeComputable_const, Complexity.polyTimeComputable_pairFstD, Complexity.polyTimeComputable_pairSndD, Randomized.BPLInP.codeStep, Randomized.BPLInP.codesW, Randomized.BPLInP.iterate_codeStep, Randomized.BPLInP.length_encWords_range_le, Turing.length_bits_le_self, Turing.length_pairEncode, Turing.pairEncode}
\difficulty{5}
For all natural numbers $B$ and $e$, the function that maps a bit string $x$ to its
associated code word is polynomial-time computable.
\end{theorem}

\begin{definition}[Binary radius word]
\lean{Randomized.BPLInP.radius}
\label{Randomized.BPLInP.radius}
\leanok
For a natural number $c$ and a Boolean word $x$, the radius word is the binary
representation of $c$ times one more than the length of the binary representation of
$|x|$.
\end{definition}

\begin{theorem}[Polynomial-Time Computability of the Radius Function]
\lean{Randomized.BPLInP.polyTimeComputable_radius}
\label{Randomized.BPLInP.polyTimeComputable_radius}
\leanok
\uses{Complexity.PolyTimeComputable, Complexity.PolyTimeComputable.append, Complexity.PolyTimeComputable.comp, Complexity.polyTimeComputable_const, Complexity.polyTimeComputable_of_linear, Complexity.polyTimeComputable_unary, Randomized.BPLInP.radius, Turing.FinTM.computesFunInTime_lengthBits}
\difficulty{5}
For every natural number $c$, the function sending a bit string $x$ to the binary
representation of $c$ times one plus the length of the binary representation of $|x|$ is
polynomial-time computable.
\end{theorem}

\begin{definition}[All-zero output word]
\lean{Randomized.BPLInP.zerosW}
\label{Randomized.BPLInP.zerosW}
\leanok
\uses{Complexity.transduce}
The map sends each Boolean word to a word of the same length consisting entirely of
false entries.
\end{definition}

\begin{theorem}[Output of the zero transducer]
\lean{Randomized.BPLInP.zerosW_eq}
\label{Randomized.BPLInP.zerosW_eq}
\leanok
\uses{Complexity.transduce, Randomized.BPLInP.zerosW}
\difficulty{4}
For every finite Boolean word $w$, the transducer that emits false on each input bit
outputs the word of length $|w|$ whose entries are all false.
\end{theorem}

\begin{definition}[Power-of-two word construction]
\lean{Randomized.BPLInP.pow2W}
\label{Randomized.BPLInP.pow2W}
\leanok
\uses{Randomized.BPLInP.zerosW}
For natural numbers $B$ and $e$ and a finite Boolean word $x$, let
$\operatorname{pow2W}(B,e,x)$ be the word consisting of $B(|x|+1)^e$ zeros followed by a
single one.
\end{definition}

\begin{theorem}[Binary value of a padded power of two]
\lean{Randomized.BPLInP.bitsVal_pow2W}
\label{Randomized.BPLInP.bitsVal_pow2W}
\leanok
\uses{Complexity.LogProg.bitsVal, Randomized.BPLInP.pow2W, Randomized.BPLInP.zerosW_eq}
\difficulty{4}
For natural numbers $B,e$ and a Boolean list $x$, the binary value of the word
consisting of $B(|x|+1)^e$ zeros followed by a one is $2^{B(|x|+1)^e}$.
\end{theorem}

\begin{theorem}[Polynomial-Time Construction of a Power-of-Two String]
\lean{Randomized.BPLInP.polyTimeComputable_pow2W}
\label{Randomized.BPLInP.polyTimeComputable_pow2W}
\leanok
\uses{Complexity.PolyTimeComputable, Complexity.PolyTimeComputable.append, Complexity.PolyTimeComputable.comp, Complexity.polyTimeComputable_const, Complexity.polyTimeComputable_polyUnary, Complexity.polyTimeComputable_transduce, Randomized.BPLInP.pow2W}
\difficulty{1}
For all natural numbers $B$ and $e$, the function sending a bit string $x$ to a string
of $B(|x|+1)^e$ zeros followed by a one is polynomial-time computable.
\end{theorem}

\begin{theorem}[Polynomial-Time Indicator for Logspace Languages]
\lean{Randomized.BPLInP.polyTimeComputable_indicator_of_LOGSPACE}
\label{Randomized.BPLInP.polyTimeComputable_indicator_of_LOGSPACE}
\leanok
\uses{Complexity.LOGSPACE, Complexity.LOGSPACE_subset_P, Complexity.PolyTimeComputable, Complexity.mem_P_iff_polyTimeComputable, Turing.MultiTapeTM.indicator}
\difficulty{1}
Let $A$ be a language over the binary alphabet that is decidable in logarithmic space.
The function that maps each bit string $z$ to the one-bit string containing the
indicator of whether $z$ belongs to $A$ is polynomial-time computable.
\end{theorem}
